\documentclass[11pt,a4paper]{article}
\usepackage[utf8]{inputenc}
\usepackage[T1]{fontenc}
\usepackage{babel}
\renewcommand{\contentsname}{Table des mati\`eres}
\usepackage{lmodern}
\usepackage[margin=2cm]{geometry}
\usepackage{xcolor}
\usepackage{titlesec}
\usepackage{longtable}
\usepackage{booktabs}
\usepackage{colortbl}
\usepackage{graphicx}
\usepackage{parskip}
\usepackage{hyperref}
\usepackage{xurl}
\usepackage{listings}
\urlstyle{tt}

\definecolor{navy}{HTML}{111111}
\definecolor{gold}{HTML}{1565C0}
\definecolor{greytext}{HTML}{555555}
\definecolor{lightgrey}{HTML}{F2F4F9}
\definecolor{okgreen}{HTML}{2E7D32}
\definecolor{warnred}{HTML}{B71C1C}

\titleformat{\section}{\Large\bfseries\color{navy}}{\thesection}{0.6em}{}[{\color{navy}\titlerule}]
\titleformat{\subsection}{\large\bfseries\color{gold}}{\thesubsection}{0.6em}{}
\titleformat{\subsubsection}{\normalsize\bfseries\color{greytext}}{\thesubsubsection}{0.6em}{}

\hypersetup{colorlinks=true, linkcolor=navy, urlcolor=navy, citecolor=navy}

\newcommand{\code}[1]{\texttt{\detokenize{#1}}}
\newcommand{\ok}{\textcolor{okgreen}{\textbf{Valid\'e en r\'eel}}}
\newcommand{\piege}{\textcolor{warnred}{\textbf{Le pi\`ege}}}
\newcommand{\avant}{\textcolor{warnred}{\textbf{Avant}}}
\newcommand{\apres}{\textcolor{okgreen}{\textbf{Apr\`es}}}

\lstdefinestyle{pythonstyle}{
  language=Python,
  basicstyle=\ttfamily\footnotesize,
  keywordstyle=\color{gold}\bfseries,
  commentstyle=\color{greytext}\itshape,
  stringstyle=\color{okgreen},
  numbers=left,
  numberstyle=\tiny\color{greytext},
  numbersep=8pt,
  frame=single,
  rulecolor=\color{lightgrey},
  backgroundcolor=\color{lightgrey},
  breaklines=true,
  showstringspaces=false,
  tabsize=2,
  xleftmargin=1em,
  framexleftmargin=0.8em,
}
\lstset{style=pythonstyle}

\lstdefinestyle{diagstyle}{
  basicstyle=\ttfamily\small\color{navy},
  backgroundcolor=\color{lightgrey},
  frame=single,
  rulecolor=\color{lightgrey},
  breaklines=false,
  xleftmargin=1em,
  framexleftmargin=0.8em,
}

\begin{document}

\begin{titlepage}
\centering
\vspace*{2cm}
{\Huge\bfseries\color{navy} Robustesse et fra\^icheur du corpus\par}
\vspace{0.5cm}
{\Large\bfseries\color{gold} Fiche p\'edagogique -- Corrections post-d\'emo\par}
\vspace{0.3cm}
{\large\color{greytext} Syst\`eme RAG pour hcp.ma\par}
\vspace{1.5cm}
{\large\color{greytext} Notes personnelles, \`a des fins d'apprentissage\par}
\vspace{0.3cm}
{\large\color{greytext} Cinq probl\`emes trouv\'es apr\`es un test en conditions r\'eelles avec l'encadrante, cinq corrections expliqu\'ees en d\'etail\par}
\vfill
\end{titlepage}

\tableofcontents
\newpage

\section{Contexte}

Un test en conditions r\'eelles avec l'encadrante a mal tourn\'e~: des r\'eponses hors
sujet, puis la m\^eme r\'eponse resservie \`a des questions diff\'erentes, dans une
session comme dans une autre. Rien de tout \c{c}a n'\'etait visible dans les 194 tests
automatis\'es qui passaient au vert -- encore une fois la m\^eme le\c{c}on d\'ej\`a tir\'ee
plusieurs fois sur ce projet~: les tests unitaires valident la \emph{logique}, seule
l'utilisation r\'eelle r\'ev\`ele ce \`a quoi la conception n'avait pas pens\'e.

La d\'emarche de diagnostic, avant toute correction~: ajouter un log temporaire
(\code{[DEBUG reformulation]}) pour voir exactement ce que le syst\`eme fait en
interne, reproduire le sc\'enario, lire la vraie base de donn\'ees et le vrai
contenu de Redis plut\^ot que de deviner. Plusieurs hypoth\`eses initiales (la
reformulation qui d\'erape, le cache qui m\'elange les sessions) se sont r\'ev\'el\'ees
fausses ou incompl\`etes en testant -- les vrais coupables \'etaient ailleurs, et il a
fallu les chercher un par un.

Cinq probl\`emes distincts ont fini par \^etre identifi\'es et corrig\'es. Cette fiche les
d\'etaille dans l'ordre o\`u ils ont \'et\'e trouv\'es.

\section{Probl\`eme 1~: un plantage silencieux sur simple panne r\'eseau}

\subsection{Ce qui a \'et\'e observ\'e}

En testant plus de questions, une erreur est apparue dans le terminal~: une
\code{ConnectionError} Python (module \code{requests}), suivie de l'arr\^et complet de
l'interface Streamlit. Aucune r\'eponse affich\'ee, toute la session perdue.

\subsection{Pourquoi}

Le module \code{Generateur} appelle Mistral pour r\'ediger la r\'eponse sur le chemin
NOTION (et pour la partie explicative du chemin MIXTE). Partout ailleurs dans le
projet, un appel \`a un LLM externe est entour\'e d'un \code{try/except} qui retombe
proprement sur un comportement par d\'efaut si l'appel \'echoue (voir le Routeur et le
Reformulateur). Ici, non~: rien ne prot\'egeait cet appel pr\'ecis. La moindre coupure
r\'eseau, le moindre hoquet de l'API Mistral, faisait remonter l'exception jusqu'en
haut de la pile et tuait tout le processus Streamlit -- pour tout le monde, pas
seulement pour la question en cours.

\begin{lstlisting}[style=diagstyle]
AVANT :
Question NOTION posee
        |
        v
Generateur.generer_reponse()
        |
        v
fonction_generation(question, contexte)  -- appel reseau vers Mistral
        |
        X  ConnectionError (reseau coupe)
        |
        v
PLANTE tout le processus Streamlit
(aucune reponse affichee, session perdue pour tous)

APRES :
Question NOTION posee
        |
        v
Generateur.generer_reponse()
        |
        v
try: fonction_generation(question, contexte)
        |
   +----+----+
   v         v
  OK      ConnectionError
   |         |
   v         v
Reponse   Message d'erreur explicite (NOTION)
normale   ou chiffre seul conserve (MIXTE)
   |         |
   +----+----+
        v
Interface reste utilisable
\end{lstlisting}

\subsection{Le correctif}

Un \code{try/except} autour de l'appel, avec une d\'egradation diff\'erente selon le
chemin~: sur NOTION, un message d'erreur explicite plut\^ot qu'un plantage. Sur MIXTE,
mieux encore~: le chiffre officiel ne d\'epend d'aucun LLM (produit par un gabarit
d\'eterministe), donc une panne de l'explication ne doit jamais priver l'utilisateur du
chiffre d\'ej\`a calcul\'e.

\begin{lstlisting}
# src/generateur.py -- chemin NOTION
texte_contexte = "\n\n".join(chunk.texte for chunk in chunks)
try:
    texte = self._fonction_generation(question, texte_contexte)
except Exception:
    return Reponse(
        texte=(
            "Une erreur technique (reseau ou service de generation "
            "indisponible) a empeche de repondre a cette question. "
            "Reessayez dans quelques instants."
        ),
        source_url="", source_titre="", source_date=None,
    )

# src/generateur.py -- chemin MIXTE
try:
    explication = self._fonction_generation(question, texte_contexte)
except Exception:
    # Le chiffre officiel reste fiable, produit sans aucun LLM --
    # une panne de l'explication ne doit pas priver l'utilisateur du chiffre.
    return Reponse(texte=phrase_chiffree, source_url=url_chiffre,
                    source_titre=titre_chiffre, source_date=date_chiffre)
\end{lstlisting}

Deux nouveaux tests reproduisent une panne r\'eseau simul\'ee (\code{ConnectionError}
lev\'ee volontairement par une fausse fonction de g\'en\'eration) sur les deux chemins~:
\ok{}.

\section{Probl\`eme 2~: un corpus fig\'e depuis des semaines}

\subsection{Ce qui a \'et\'e observ\'e}

L'encadrante a interrog\'e le syst\`eme sur une publication r\'ecente~: rien trouv\'e, ou
une r\'eponse construite \`a partir d'un contenu p\'erim\'e. En v\'erifiant directement la
base de donn\'ees, le document le plus r\'ecemment publi\'e qui soit index\'e datait de
plusieurs semaines -- rien de plus r\'ecent n'existait dans le syst\`eme.

\subsection{Pourquoi}

L'automatisation de la mise \`a jour \'etait pr\'evue d\`es le d\'ebut du projet (voir
l'ADR sur l'API BDS~: \emph{"un script planifi\'e, ex. une fois par jour"}), mais cette
\'etape n'avait jamais \'et\'e r\'eellement c\^abl\'ee \`a un vrai planificateur -- ni t\^ache
Windows, ni cron. Les scripts existaient (\code{indexer_documents.py},
\code{preremplir_indicateurs_bds.py}), mais rien ne les d\'eclenchait tout seul~: ils
devaient \^etre relanc\'es \`a la main, et \c{c}a n'avait pas \'et\'e fait depuis longtemps.

\subsection{Le correctif}

Un script d'orchestration Python enchaine les deux \'etapes (voir le code plus bas),
appel\'e chaque nuit par une v\'eritable t\^ache planifi\'ee Windows, via un petit script
PowerShell interm\'ediaire.

\begin{lstlisting}[style=diagstyle]
Planificateur de taches Windows (chaque nuit)
        |
        v
rafraichir_corpus.ps1
        |
        v
python -m scripts.rafraichir_corpus
        |
        +--> indexer_documents.main()
        |       parcourt la page 1 de chaque categorie hcp.ma
        |
        +--> preremplir_indicateurs_bds.main()
        |       met a jour les 18 indicateurs BDS cures
        |
        +--> (voir Probleme 4 -- decide s'il faut vider le cache)
\end{lstlisting}

\begin{lstlisting}
schtasks /create /tn "RAG_HCP_Rafraichissement_Nocturne" ^
  /tr "powershell.exe -ExecutionPolicy Bypass -File rafraichir_corpus.ps1" ^
  /sc daily /st 03:00
\end{lstlisting}

Limite assum\'ee~: cette t\^ache ne s'ex\'ecute que si le PC est allum\'e \`a l'heure
pr\'evue -- un palliatif pour la phase de d\'eveloppement sur machine personnelle. Une
fois d\'eploy\'e sur un vrai serveur (24/7 par d\'efinition), un cron classique remplace
directement le Planificateur Windows, sans changer une ligne du code Python.

Test\'e en conditions r\'eelles~: le corpus est pass\'e de 96 \`a 99 documents (7499
chunks) d\`es le premier rafra\^ichissement -- \ok{}.

\section{Probl\`eme 3~: relancer l'indexation la polluait}

\subsection{Ce qui a \'et\'e observ\'e}

En pr\'eparant le correctif pr\'ec\'edent, un examen du code d'\code{indexer_documents.py}
a r\'ev\'el\'e un second bug, jamais rencontr\'e jusque-l\`a parce que ce script n'avait
jamais \'et\'e relanc\'e plusieurs fois sur le m\^eme corpus.

\subsection{Pourquoi}

La page 1 de chaque listing hcp.ma contient toujours un m\'elange de publications
d\'ej\`a connues et de vraies nouveaut\'es (le site trie du plus r\'ecent au plus ancien,
mais une page enti\`ere d\'epasse presque toujours ce qui a chang\'e depuis la derni\`ere
fois). La contrainte \code{document.url} UNIQUE emp\^echait bien un doublon dans la
table \code{document}, mais rien n'emp\^echait de refaire extraction, d\'ecoupage et
calcul des embeddings pour un document \emph{d\'ej\`a index\'e}, et de r\'einjecter des
chunks en double dans la base et l'index -- co\^uteux, et \c{c}a polluait le corpus \`a
chaque relance.

\begin{lstlisting}[style=diagstyle]
AVANT (a chaque relance) :          APRES :
Page 1 rencontree                   Page 1 rencontree
   |                                   |
   v                                   v
Deja connue ? oui                  Deja connue ? oui
   |                                   |
   v                                   v
Reextraction + reembedding         ARRET immediat, rien de plus
+ chunks dupliques en base         (statut "deja_indexe")
\end{lstlisting}

\subsection{Le correctif}

Une v\'erification pr\'ecoce, avant tout travail co\^uteux~: si l'URL est d\'ej\`a en base,
on s'arr\^ete net.

\begin{lstlisting}
# scripts/indexer_documents.py
deja_connu = conn.execute(
    "SELECT id_document FROM document WHERE url = ?", (document.url,)
).fetchone()
if deja_connu is not None:
    return {"statut": "deja_indexe", "chunks": 0, "tableaux": 0,
            "id_document": deja_connu[0]}

texte, tableaux = extracteur.extraire(document)  # seulement si vraiment nouveau
\end{lstlisting}

Un test existant v\'erifiait d\'ej\`a qu'une m\^eme URL ne cr\'eait pas deux lignes
\code{document} -- mais aucun test ne v\'erifiait le nombre de \emph{chunks}, ce qui
avait laiss\'e ce bug invisible. Un nouveau test comble ce trou pr\'ecis~: \ok{}.

\section{Probl\`eme 4~: le cache peut resservir une r\'eponse p\'erim\'ee}

\subsection{Le dilemme}

Le cache des r\'eponses NOTION (Redis) est volontairement partag\'e entre toutes les
sessions~: si une personne demande "pourquoi le ch\^omage a-t-il augment\'e~?"
aujourd'hui, une autre personne qui pose la m\^eme question demain doit pouvoir
b\'en\'eficier d'une r\'eponse d\'ej\`a calcul\'ee, sans rappeler le LLM inutilement -- c'est
le but assum\'e du cache, pas un oubli \`a corriger.

Mais un TTL de 24h glissant depuis l'\'ecriture de l'entr\'ee n'est pas align\'e sur le
moment r\'eel o\`u le corpus change~: le rafra\^ichissement nocturne tourne \`a heure fixe
(3h), alors qu'une entr\'ee mise en cache la veille \`a 14h reste valide jusqu'\`a 14h le
lendemain -- jusqu'\`a 11h de d\'ecalage o\`u le cache peut resservir une r\'eponse plus
vieille que du contenu d\'ej\`a mis \`a jour.

\subsection{Options \'ecart\'ees, et pourquoi}

Cloisonner le cache par session (une cl\'e Redis par \code{id_session}) r\'eglait le
probl\`eme de fra\^icheur, mais au prix du b\'en\'efice de partage entre sessions --
\'ecart\'e, ce b\'en\'efice \'etant justement ce qui \'etait recherch\'e d\`es le d\'epart.

Une invalidation cibl\'ee (ne supprimer que les entr\'ees concern\'ees par un nouveau
document) demanderait de savoir, pour chaque r\'eponse en cache, quels documents ont
servi \`a la produire -- puis de comparer s\'emantiquement chaque nouveau document \`a
chaque question en cache, en utilisant la m\^eme machinerie que la recherche
(embeddings, d\'ecoupage, fusion RRF). Cette machinerie existe d\'ej\`a dans le projet,
mais elle ne tourne qu'au moment o\`u un utilisateur pose une vraie question -- la
d\'eclencher pour comparer chaque nouveau document \`a tout le cache serait co\^uteux, et
pas fiable \`a 100\,\% de toute fa\c{c}on. \'Ecart\'e.

\subsection{La d\'ecision retenue}

Garder le cache global, mais le vider enti\`erement -- sans distinction de sujet --
d\`es qu'au moins une vraie nouveaut\'e est trouv\'ee pendant le rafra\^ichissement
nocturne. Simple, fiable \`a 100\,\%, et n'affecte pas les nuits o\`u rien n'a chang\'e.

\begin{lstlisting}[style=diagstyle]
Rafraichissement nocturne termine
        |
        v
Au moins 1 nouveau document trouve ?
        |
   +----+----+
   v         v
  oui        non
   |         |
   v         v
Cache NOTION   Cache NOTION
vide entier    conserve tel quel
(toutes cles)  (rien ne change)
\end{lstlisting}

\begin{lstlisting}
# scripts/vider_cache_notion.py
def vider(client_redis) -> int:
    if client_redis is None:
        return 0
    cles = list(client_redis.scan_iter(PREFIXE_CACHE + "*"))
    if cles:
        client_redis.delete(*cles)
    return len(cles)

# scripts/rafraichir_corpus.py -- decision
nouveaux_documents = indexer_documents.main(chemin_db=chemin_db)
...
if nouveaux_documents > 0:
    client = _connecter_redis()
    vider(client)
else:
    print("Aucune nouveaute -> cache NOTION conserve.")
\end{lstlisting}

Limite assum\'ee, discut\'ee explicitement~: hcp.ma ne semble jamais r\'eviser un
document existant sous la m\^eme URL, seulement publier de nouveaux documents avec des
chiffres \`a jour. Un nouveau document peut donc rendre une r\'eponse en cache obsol\`ete
sur un sujet qu'il ne partage avec \emph{aucun} document d\'ej\`a cit\'e dans le cache --
un suivi "quel document a servi \`a quelle r\'eponse" ne suffirait pas \`a lui seul \`a
cibler quoi invalider. Vider tout d\`es qu'il y a une nouveaut\'e reste la seule
garantie fiable sans construire la comparaison s\'emantique compl\`ete \'evoqu\'ee
plus haut.

4 nouveaux tests (suppression effective, historique jamais touch\'e, cache vide ou
absent g\'er\'e proprement)~: \ok{}.

\section{Probl\`eme 5~: un historique qui ne s'effa\c{c}ait jamais}

\subsection{Ce qui a \'et\'e observ\'e}

En inspectant directement Redis (\code{redis-cli}, ou l'\'equivalent en Python), les
cl\'es d'historique de conversation affichaient toutes \code{TTL: -1} -- aucune
expiration. Elles s'accumulaient ind\'efiniment, contrairement au cache qui, lui, avait
toujours eu 24h.

\subsection{Le correctif}

Un TTL glissant de 24h, repouss\'e \`a chaque nouveau message de la session plut\^ot que
fix\'e une seule fois \`a la cr\'eation -- une conversation active ne doit jamais
expirer en plein milieu, seule une session vraiment abandonn\'ee finit par dispara\^itre.

\begin{lstlisting}[style=diagstyle]
Session active :
  message 1 -> rpush + expire(24h)       TTL = 24h
  message 2 -> rpush + expire(24h)       TTL reparti a 24h (pas de coupure)
  ...

Session abandonnee (onglet ferme sans rien faire d'autre) :
  aucun nouveau message -> TTL continue de descendre -> 0 -> cle supprimee
\end{lstlisting}

\begin{lstlisting}
# src/cache_redis.py
def ajouter(self, id_session: str, question: str, reponse: Reponse) -> None:
    if self._client is None:
        return
    entree = json.dumps({"question": question, "reponse": asdict(reponse)})
    cle = PREFIXE_HISTORIQUE + id_session
    self._client.rpush(cle, entree)
    self._client.expire(cle, TTL_HISTORIQUE)  # 24h, repousse a chaque ajout
\end{lstlisting}

Une alternative envisag\'ee -- supprimer l'historique \`a la fermeture exacte de
l'onglet -- a \'et\'e \'ecart\'ee~: Streamlit n'offre aucun signal fiable c\^ot\'e serveur
pour d\'etecter ce moment pr\'ecis (pas de vraie notification de d\'econnexion
exploitable simplement). Un bouton "Nouvelle conversation" explicite reste possible
plus tard pour le cas o\`u l'utilisateur veut clore volontairement -- pas construit
pour l'instant, le TTL glissant \'etant jug\'e suffisant.

2 nouveaux tests (TTL pos\'e \`a la cr\'eation, TTL repouss\'e \`a chaque nouvel ajout)~:
\ok{}.

\section{Une d\'ecision \`a part~: garder le log de diagnostic}

Le log \code{[DEBUG reformulation]} ajout\'e pour cette enqu\^ete (affiche dans le
terminal, \`a chaque reformulation, le texte avant et apr\`es) a failli \^etre retir\'e
une fois le diagnostic termin\'e -- puis conserv\'e volontairement~: il r\'epond
exactement \`a un besoin de visibilit\'e d\'ej\`a demand\'e pour cette phase de
d\'eveloppement, sans qu'il ait \'et\'e \'evident au d\'epart que c'\'etait techniquement
possible aussi simplement.

\begin{lstlisting}
# scripts/poser_question.py
question_effective = reformuler_si_necessaire(question, entrees_recentes, fonction_reformulation)
if question_effective != question:
    print(f"[DEBUG reformulation] '{question}' -> '{question_effective}'")
type_question = routeur.classifier(question_effective)
\end{lstlisting}

\section{Les tests~: ce qui a \'et\'e valid\'e}

\begin{longtable}{@{}p{6.8cm}p{2cm}p{6.7cm}@{}}
\rowcolor{navy}
\textcolor{white}{\textbf{Correctif}} & \textcolor{white}{\textbf{Tests}} &
\textcolor{white}{\textbf{Ce qui est prouv\'e}} \\
\midrule
\endhead
D\'egradation propre du \code{Generateur} & 2 &
Panne r\'eseau simul\'ee sur NOTION (message d'erreur, pas de plantage) et sur MIXTE
(chiffre conserv\'e). \\
Anti-duplication \code{indexer_documents} & 1 &
Une m\^eme URL relanc\'ee deux fois ne cr\'ee plus de chunks en double. \\
\code{vider_cache_notion} & 4 &
Suppression effective du cache, historique jamais touch\'e, cas sans client g\'er\'e
proprement. \\
TTL glissant sur l'historique & 2 &
TTL de 24h pos\'e \`a la cr\'eation, repouss\'e \`a chaque nouveau message. \\
\end{longtable}

Suite compl\`ete du projet~: \textbf{194 tests avant ces corrections} $\rightarrow$
\textbf{196 tests d\'esormais}, tous verts. Aucune r\'egression sur l'existant \`a
aucune \'etape.

\section{Synth\`ese -- le fil conducteur}

\textbf{Un test r\'eel avec un tiers vaut plus que des semaines de tests solo.} Les
cinq probl\`emes de cette fiche existaient depuis un moment, invisibles en testant
seul avec les m\^emes questions habituelles -- il a fallu quelqu'un d'ext\'erieur, avec
des questions impr\'evues, pour les faire remonter d'un coup.

\textbf{Diagnostiquer avant de corriger, toujours.} Plusieurs pistes initiales
(reformulation qui d\'erape, cache qui m\'elange les sessions) se sont r\'ev\'el\'ees
fausses ou incompl\`etes une fois v\'erifi\'ees avec un vrai log et une vraie inspection
de Redis. Corriger \`a l'aveugle sur une hypoth\`ese non v\'erifi\'ee aurait pu masquer
le vrai probl\`eme au lieu de le r\'esoudre.

\textbf{La d\'egradation propre doit \^etre syst\'ematique, pas au cas par cas.} Le
Routeur et le Reformulateur avaient d\'ej\`a ce r\'eflexe (un \code{try/except} autour de
chaque appel externe) -- le \code{Generateur} \'etait le seul endroit oubli\'e, et \c{c}a a
suffi \`a faire planter toute l'interface. Une r\`egle de conception n'a de valeur que
si elle est appliqu\'ee partout, sans exception oubli\'ee.

\textbf{Simple et fiable bat sophistiqu\'e et fragile.} Vider tout le cache d\`es qu'une
nouveaut\'e appara\^it est moins \'el\'egant qu'une invalidation cibl\'ee par sujet -- mais
cette derni\`ere demanderait une comparaison s\'emantique co\^uteuse et jamais garantie
\`a 100\,\%. \`A l'\'echelle de ce projet, la solution simple est aussi la plus s\^ure.

\section{Ce qui reste ouvert}

\begin{itemize}
  \item \textbf{D\'ecouverte de nouveaut\'es limit\'ee \`a la page 1} de chaque cat\'egorie
    -- coh\'erent avec l'objectif "fra\^icheur quotidienne", mais \`a surveiller si
    hcp.ma publie un jour plusieurs pages de nouveaut\'es en une seule nuit
    (rafra\^ichissement manqu\'e, par exemple).
  \item \textbf{Couverture des pages de listing} pas encore v\'erifi\'ee manuellement
    contre le vrai site -- \`a confirmer qu'aucune cat\'egorie ou sous-section active
    n'est absente de \code{data/listing_urls.py}.
  \item \textbf{Bouton "Nouvelle conversation"} évoqu\'e mais pas construit -- le TTL
    glissant de 24h a \'et\'e jug\'e suffisant pour l'instant.
  \item \textbf{Jeu de test de fiabilit\'e} toujours \`a construire (30 \`a 50
    questions-r\'eponses, mesure du taux d'exactitude) -- prochaine \'etape, une fois
    ces corrections consid\'er\'ees stables.
\end{itemize}

\end{document}
